\section{Local physical ranges of G-injective tensors}%
\label{sec:peps_g_injective_physical_ranges}

\begin{theorem}[Invariant virtual space and physical range]%
    \label{thm:peps_g_injective_invariants_range}
    \lean{TNLean.PEPS.IsGInjective.ker_eq_ker_averageMap,
        TNLean.PEPS.IsGInjective.bijective_invariants_rangeRestrict,
        TNLean.PEPS.IsGInjective.invariantsRangeEquiv,
        TNLean.PEPS.IsGInjective.finrank_range}
    \leanok
    \uses{def:peps_g_injective}
    Let $G$ be a finite group acting linearly on a complex virtual space
    $\mathcal V$, and let $P_G=|G|^{-1}\sum_{g\in G}U_g$ be the
    projector onto $\mathcal V^G$. If
    $T:\mathcal V\to\mathcal P$ is $G$-injective, then
    \[
        \ker T=\ker P_G,
        \qquad T\big|_{\mathcal V^G}:\mathcal V^G
        \mathrel{\cong}\operatorname{im}T.
    \]
    In finite dimension, $\dim\operatorname{im}T=\dim\mathcal V^G$.
    These are local consequences of
    \cite[Definition~5.1]{Schuch2010PEPS}.
    The virtual projector can be inserted before the physical map:
    % Formula: T P_G = T, the invariant-boundary form of Definition 5.1(i).
    % Source: Papers/1001.3807/paper_v3.tex:1278--1296, figs3/ug-sym.
    % Ink-to-index: a west virtual leg is the entire boundary configuration
    % in V, not one microscopic PEPS bond; the east physical leg is in P.
    % Contraction: the P_G--T join sums over a basis of V. Both panels have
    % one open virtual input and one open physical output, hence (V,P)=(1,1).
    \begin{tenkzequation}
        \tenkzkernel
        \begin{tenkz}[rows={wire}, cols=2, bonds=none]
            \tn[name=rangeProjector, skin=box,
                ports={180:virtual, 0:virtual}]{P_G}
            \tn[at=(1,2), name=rangeTensor,
                ports={180:virtual, 0:physical}]{T}
            \tnwire{rangeProjector.0}{rangeTensor.180}
        \end{tenkz}
        $=$
        \begin{tenkz}[rows={wire}, cols=1, bonds=none]
            \tn[ports={180:virtual, 0:physical}]{T}
        \end{tenkz}
    \end{tenkzequation}
\end{theorem}

\begin{proof}\leanok
    Invariance gives $T P_G=T$. If $Tx=0$, then $T P_Gx=0$;
    injectivity on $\mathcal V^G$ implies $P_Gx=0$. The converse
    follows from $T P_G=T$. The restriction is injective by
    $G$-injectivity, and every $Tx$ equals $T(P_Gx)$, proving
    surjectivity onto the physical range.
\end{proof}

\begin{theorem}[Comparison of physical ranges]%
    \label{thm:peps_g_injective_physical_range_equivalence}
    \lean{TNLean.PEPS.IsGInjective.rangeEquiv,
        TNLean.PEPS.IsGInjective.rangeEquiv_apply,
        TNLean.PEPS.IsGInjective.rangeEquiv_unique}
    \leanok
    \uses{thm:peps_g_injective_invariants_range}
    Let $T:\mathcal V\to\mathcal P$ and
    $S:\mathcal V\to\mathcal Q$ be $G$-injective for the same
    virtual representation. There is a unique linear map
    $F:\operatorname{im}T\to\operatorname{im}S$ such that
    $F(Tx)=Sx$ for every $x\in\mathcal V$, and this map is an
    isomorphism.
\end{theorem}

\begin{proof}\leanok
    Compose the inverse of
    $T|_{\mathcal V^G}$ with $S|_{\mathcal V^G}$.
    Since $Tx=T(P_Gx)$ and $Sx=S(P_Gx)$, the resulting map has the
    required identity. That identity specifies its value at every
    element of $\operatorname{im}T$, so it also proves uniqueness.
\end{proof}
